\documentclass[11pt]{article}
\usepackage[margin=0.82in]{geometry}
\usepackage{amsmath,amssymb,amsthm,mathtools,microtype}
\usepackage[T1]{fontenc}
\usepackage{lmodern}
\usepackage{graphicx}
\usepackage[hidelinks]{hyperref}
\newtheorem{theorem}{Theorem}
\newtheorem{corollary}{Corollary}
\newtheorem{lemma}{Lemma}
\newcommand{\R}{\mathbb R}
\newcommand{\N}{\mathbb N}
\newcommand{\V}{\mathcal V}
\newcommand{\supp}{\operatorname{supp}}
\newcommand{\Ran}{\operatorname{Ran}}
\newcommand{\Dom}{\operatorname{Dom}}
\setlength{\parindent}{0pt}
\setlength{\parskip}{4pt}
\title{Exact transmission conditions and inclusions\\
for finite-step variable bandwidth}
\author{Automated mathematical discovery run \texttt{fa\_banach\_001}, lane 11}
\date{17 August 2026}
\begin{document}
\maketitle

\textbf{Status.} Likely valid full answers to Questions II.5 and II.6 of
arXiv:2503.22443, plus a sharp one-sided result toward the sampled-sign
question.

\section*{Source questions and normalization}

Jaming and Perez ask when local Paley--Wiener pieces patch to their
finite-step variable-bandwidth space, and whether an order between two step
coefficients implies inclusion of the corresponding spaces.  The displayed
comparison in Question II.6 reads $p_1\ge p_2$ although the two functions are
called $p,\widetilde p$; below it is read in the natural intended sense
$p\ge\widetilde p$ pointwise.

\begin{center}
\includegraphics[width=.61\textwidth]{figures/open_questions_crop.png}
\end{center}

It is cleaner to separate the spectral cutoff from the paper's competing
$\Omega$ normalizations.  Let
\[
 A_p f=-(pf')',\qquad
 \V_p(B)=\Ran\mathbf 1_{[0,B]}(A_p),\qquad B>0.
\]
Suppose $p=p_j>0$ on $I_j=(x_j,x_{j+1})$, where
$-\infty=x_0<x_1<\cdots<x_N<x_{N+1}=\infty$, and put
\[
                         c_j=\sqrt{B/p_j}.
\]
Thus an element of $\V_p(B)$ restricts on $I_j$ to an entire member of the
classical $PW_{c_j}(\R)$.  In the explicit formulas of the source,
$c_j=\Omega p_j^{-1/2}$, so $B=\Omega^2$.

\section*{Full answer to the local-patching question}

For each $j$, let $F_j\in PW_{c_j}(\R)$, and define
$f(x)=F_j(x)$ for $x\in I_j$.  Values and derivatives below are those of the
entire functions $F_j$.

\begin{theorem}[Transmission-jet characterization]\label{thm:jets}
The piecewise function $f$ belongs to $\V_p(B)$ if and only if, at every jump
$x_j$ and for every $m\ge0$,
\begin{align}
 p_{j-1}^{m}F_{j-1}^{(2m)}(x_j)
   &=p_j^{m}F_j^{(2m)}(x_j),                         \tag{T$_{2m}$}\label{eq:even}\\
 p_{j-1}^{m+1}F_{j-1}^{(2m+1)}(x_j)
   &=p_j^{m+1}F_j^{(2m+1)}(x_j).                    \tag{T$_{2m+1}$}\label{eq:odd}
\end{align}
In particular, continuity is only the first of an infinite, explicit family
of matching conditions.
\end{theorem}

\begin{proof}
If $f\in\V_p(B)$, then $f\in\Dom(A_p^m)$ for every $m$.  A function $u$ in
$\Dom(A_p)$ is continuous at every jump and has continuous flux $pu'$ there.
On $I_j$,
\[
                  A_p^m f=(-1)^m p_j^mF_j^{(2m)}.    \tag{1}
\]
Applying continuity and flux continuity to $A_p^m f$ gives
\eqref{eq:even} and \eqref{eq:odd}.

Conversely, the conditions with $m=0$ put $f$ in $\Dom(A_p)$.  The complete
family, used inductively, puts $f$ in every $\Dom(A_p^m)$ and proves (1).
Bernstein's inequality in $PW_{c_j}$ now gives, with
$C^2=\sum_j\|F_j\|_2^2$,
\begin{align*}
 \|A_p^m f\|_2^2
 &\le \sum_j p_j^{2m}\|F_j^{(2m)}\|_2^2
 \le \sum_j p_j^{2m}c_j^{4m}\|F_j\|_2^2\\
 &= C^2B^{2m}.                                      \tag{2}
\end{align*}
Let $\mu_f$ be the spectral measure of $f$ for $A_p$.  If
$\mu_f((B+\varepsilon,\infty))=\delta>0$, then the spectral theorem gives
\[
 \|A_p^m f\|_2^2=\int\lambda^{2m}\,d\mu_f(\lambda)
 \ge\delta(B+\varepsilon)^{2m},
\]
contradicting (2) as $m\to\infty$.  Hence $\supp\mu_f\subset[0,B]$ and
$f\in\V_p(B)$.
\end{proof}

The conditions also reveal why one local piece determines the global
function.  At a jump $x$, write the adjacent constants as $p_-,p_+$ and set
$H_-(t)=F_-(x+\sqrt{p_-}\,t)$.  Then the unique compatible right piece is
given by
\begin{equation}
\begin{aligned}
 H_+(t):=F_+(x+\sqrt{p_+}\,t)
 &=\frac{H_-(t)+H_-(-t)}2 \\
 &\quad+\sqrt{\frac{p_-}{p_+}}
       \frac{H_-(t)-H_-(-t)}2.
\end{aligned}
\tag{3}\label{eq:transfer}
\end{equation}
Indeed, the even Taylor coefficients in (3) encode \eqref{eq:even}, and the
odd ones encode \eqref{eq:odd}.  Both sides have exponential type at most
$\sqrt B$ in $t$, and reflection and dilation preserve $L^2$ membership.
Iteration and Theorem~\ref{thm:jets} yield the following useful strengthening
of the source's uniqueness observation.

\begin{corollary}[Onto local chart]\label{cor:chart}
For every component $I_j$, restriction followed by entire continuation is a
real-linear bijection
\[
            \V_p(B)_{\R}\longrightarrow PW_{\sqrt{B/p_j}}(\R)_{\R}.
\]
\end{corollary}

\section*{Full answer to the inclusion question}

The next theorem allows the finite jump sets of $p$ and $\widetilde p$ to be
different.

\begin{theorem}[Complete inclusion classification]\label{thm:inclusion}
For positive finite-step functions $p,\widetilde p$ and $B>0$,
\[
 \V_p(B)\subseteq\V_{\widetilde p}(B)
 \quad\Longleftrightarrow\quad
 \widetilde p=\alpha p\ \hbox{ for some }0<\alpha\le1.              \tag{4}
\]
Consequently, even under the pointwise order $p\ge\widetilde p$, inclusion
fails unless the coefficients are global scalar multiples.
\end{theorem}

\begin{proof}
Assume the inclusion.  At any breakpoint $x$ in the union of the two step
partitions, denote the one-sided constants by $p_\pm,\widetilde p_\pm$;
the two constants are equal if the corresponding function has no jump there.
Every $f\in\V_p(B)$ must satisfy both flux laws
\[
 p_-f'(x-)=p_+f'(x+),\qquad
 \widetilde p_-f'(x-)=\widetilde p_+f'(x+).            \tag{5}
\]
If $p_-/p_+\ne\widetilde p_-/\widetilde p_+$, the two equations force both
derivatives to vanish for every $f\in\V_p(B)$.  This is impossible by
Corollary~\ref{cor:chart}: in either adjacent local chart, prescribe a real
Paley--Wiener function whose derivative at $x$ is nonzero.  Thus the ratios
agree at every breakpoint.  The quotient $\widetilde p/p$ is consequently
constant across the common refinement, say $\widetilde p=\alpha p$.

Conversely, $A_{\alpha p}=\alpha A_p$ on the same domain, and therefore
\[
                      \V_{\alpha p}(B)=\V_p(B/\alpha). \tag{6}
\]
Equation (6) gives the inclusion exactly when $\alpha\le1$.  It is strict in
the opposite direction: a local chart contains functions with Fourier mass
between the two distinct local band edges.
\end{proof}

Here is a concrete negative answer to the proposed monotonicity.  With one
jump at zero, take
\[
 p(x)=\begin{cases}4,&x<0,\\1,&x>0,\end{cases}\qquad
 \widetilde p(x)=\begin{cases}2,&x<0,\\1,&x>0.\end{cases}            \tag{7}
\]
Then $p\ge\widetilde p$.  In the left local chart let
$c=\sqrt B/2$ and prescribe
\[
             \widehat F(\xi)=-i\xi\,\mathbf1_{[-c,c]}(\xi).
\]
This defines a real $F\in PW_c$ with
$F'(0)=2c^3/(3\sqrt{2\pi})>0$.  Its lift lies in $\V_p(B)$.  If it also lay
in $\V_{\widetilde p}(B)$, (5) would simultaneously give
$f'(0+)=4f'(0-)$ and $f'(0+)=2f'(0-)$, a contradiction.

\section*{Sharp tail sampling and an all-bandwidth set}

The source's final question asks for discrete sign-retrieval sets.  The full
multi-component classification remains open, but the local chart gives a
sharp answer for a natural family.

\begin{theorem}[Sharp one-sided lattice]\label{thm:sampling}
Fix the right unbounded component $I_N=(x_N,\infty)$, choose $a>x_N$, and put
$c=\sqrt{B/p_N}$.  The one-sided lattice
\[
                         \Lambda_h=\{a+nh:n\ge0\}
\]
sign retrieves every real member of $\V_p(B)$ if and only if
\[
                              0<h<\frac{\pi}{2c}.       \tag{8}
\]
\end{theorem}

\begin{proof}
By Corollary~\ref{cor:chart}, it suffices to work in real $PW_c$.  Suppose
$h<\pi/(2c)$ and $|F|=|G|$ on $\Lambda_h$.  Then
$H=F^2-G^2\in PW_{2c}$ and $H(a+nh)=0$ for every $n\ge0$.  If
$\varphi(\xi)=e^{ia\xi}\widehat H(\xi)$, these zeros say that $\varphi$ is
orthogonal on $[-2c,2c]$ to every nonnegative mode $e^{inh\xi}$.
After the change of variable $\theta=h\xi$, extend $\varphi(\theta/h)$ by
zero to the circle.  Its support is a proper arc because $4ch<2\pi$, while
the mode conditions put it in one Hardy boundary subspace.  Hardy boundary
uniqueness forces it to vanish: a nonzero Hardy function cannot have boundary
values zero on the complementary arc of positive measure.  Thus $H=0$.
Since entire functions form an integral domain,
$(F-G)(F+G)=0$ implies $F=G$ or $F=-G$.

For sharpness, let $h\ge\pi/(2c)$ and $d=\pi/(2h)\le c$.  Fill the removable
singularities in
\[
 U(z)=\frac{\sin(d(z-a))}{z-a+2h},\qquad
 V(z)=\frac{\cos(d(z-a))}{z-a+h}.                      \tag{9}
\]
These are nonzero real members of $PW_c$: they have type $d$ and are square
integrable on $\R$.  At even samples $U=0$, and at odd samples $V=0$.
Hence $F=U+V$ and $G=U-V$ have the same sampled magnitudes but are not global
signs of one another.  Their local-chart lifts prove failure in $\V_p(B)$.
\end{proof}

There is also one discrete set that works for every finite bandwidth at once.
For $a>x_N$, let
\[
                         \Lambda_\infty=\{a+\sqrt n:n\ge1\}.          \tag{10}
\]
If a nonzero entire function has finite exponential type, Jensen's formula
gives only $O(r)$ zeros in $|z|\le r$, whereas (10) contributes order $r^2$.
Thus (10) is a uniqueness set for $PW_{2c}$ for every finite $c$.  Applying
the squared-function argument above proves simultaneous sign retrieval in
$\V_p(B)_\R$ for every $B>0$.  This resolves the literal bandwidth-free
reading of the source's notation $PW(A_p)$, though the set is not separated
and no stability claim is made.

\section*{Scope, novelty, and human-review focus}

Theorems~\ref{thm:jets} and \ref{thm:inclusion} fully answer the two questions
reproduced in the crop.  Theorem~\ref{thm:sampling} and (10) answer substantial
one-component cases of the final sampled-sign question, not its arbitrary
multi-component classification.

A bounded search through 17 August 2026 used the exact source title and
questions together with ``variable bandwidth inclusion,'' ``step
Sturm--Liouville spectral subspace,'' and ``sampled sign retrieval.''  It
found the source and Thakur's classical unsigned Paley--Wiener theorem, but no
answer to these finite-step questions.  Novelty is therefore only
``apparently new within the bounded search.''

Human review should focus on (i) the source's inconsistent $\Omega$ cutoff,
which is why the packet uses the unambiguous edge $B$; (ii) the induction from
the jet conditions to $\Dom(A_p^m)$; and (iii) the intended reading of the
typo in Question II.6.  None affects the operator-theoretic statement in the
$B$ normalization.

\section*{References}

\begin{enumerate}
\item P. Jaming and R. Perez III, ``Sign retrieval in spaces of variable
  bandwidth,'' arXiv:2503.22443 (2025), Questions II.5, II.6, and Section IV.
\item K. Gr\"ochenig and A. Klotz, ``What is variable bandwidth?''
  Comm. Pure Appl. Math. 70 (2017), 2039--2083.
\item G. Thakur, ``Reconstruction of bandlimited functions from unsigned
  samples,'' J. Fourier Anal. Appl. 17 (2011), 720--732;
  arXiv:1007.0198.
\end{enumerate}

\end{document}
